\documentclass[12pt]{report}
\usepackage{amsmath}
\usepackage{amssymb}
\usepackage{amsthm}
\usepackage{enumitem}
\usepackage{multicol}

\begin{document}
\renewcommand\qedsymbol{OE$\Delta$}
\newcommand{\abs}[1]{\left| #1 \right|}
\newcommand{\mathif}{\text{if }}
\newcommand{\seq}[2][n]{\ensuremath{(#2 _{#1})}}

\tableofcontents
\pagebreak

\setcounter{chapter}{2}
\chapter{A Game of Hide and Sequence}

\section{Exercise 3.1}
    Suppose that a sequence ($a_n$) converges to 0.001. 
    Prove that finitely many values of $a_n$ are negative. 

    \subsection{Solution}
        \begin{proof}
            Recall the definition of convergence.
            \begin{equation}
                \forall \varepsilon > 0, \exists N : \forall n > N, \left| a_n - 0.001 \right| < \varepsilon
            \end{equation}

            The last part of that can have the absolute value removed.
            \begin{equation}
                -\varepsilon < a_n - 0.001 < \varepsilon
            \end{equation}

            Suppose that $\varepsilon = 0.001$.
            \begin{gather}
                -0.001 < a_n - 0.001 < 0.001\\
                0 < a_n < 0.002
            \end{gather}

            This means that for all $n > N$, $a_n$ is positive.
            This means that a finite number $N$ of $a_n$ can be negative or zero.
            This means that a finite number of $a_n$ are negative.
        \end{proof}    
    \pagebreak
    
\section{Exercise 3.22}
    Suppose ($a_n$) and ($b_n$) are sequences.
    \begin{enumerate}[label=(\alph*)]
        \item   Prove that $a_n \to L$ and $a_n \leq M$ for all $n$, then $L \leq M$.
        \item   Assume that $a_n \leq b_n$ for all $n$. Prove that if $a_n \to L$ and $b_n \to M$, then $L \leq M$.
    \end{enumerate}

    \subsection{Solution (b)}
        \begin{proof}[Proof by Contradiction]
            Suppose for contradiction that $L > M$.

            We know from the start that $a_n \to L$ and $b_n \to M$.
            \begin{gather}
                \label{eq:3.22.b.1}
                \forall \varepsilon > 0, \exists N : \forall n > N, \left| a_n - L \right| < \varepsilon\\
                \label{eq:3.22.b.2}
                \forall \varepsilon > 0, \exists N : \forall n > N, \left| b_n - M \right| < \varepsilon
            \end{gather}

            Furthermore, we can subtract $b_n$ from $a_n$, which also has a limit.
            \begin{gather}
                \label{eq:3.22.b.3}
                b_n - a_n \to M - L\\
                \label{eq:3.22.b.4}
                \forall \varepsilon > 0, \exists N : \forall n > N, \left| (b_n - a_n) - (M - L) \right| < \varepsilon
            \end{gather}

            In both (\ref{eq:3.22.b.1}) and (\ref{eq:3.22.b.2}), $\varepsilon > 0$.
            We can recall $L > M$, just so it's enumerated.
            \begin{gather}
                \label{eq:3.22.b.5}
                L > M\\
                \label{eq:3.22.b.6}
                L - M > 0\\
                \label{eq:3.22.b.7}
                \left| L - M \right| = L - M
            \end{gather}

            Take (\ref{eq:3.22.b.2}) and manipulate it into a similar equation by multiplying by $-1$ in the absolute value (it changes nothing else).
            \begin{equation}
                \label{eq:3.22.b.8}
                \left| M - b_n \right| < \varepsilon
            \end{equation}

            We can add together (\ref{eq:3.22.b.1}) and (\ref{eq:3.22.b.8}) and get another $\varepsilon > 0$.
            \begin{equation}
                \label{eq:3.22.b.9}
                \left| a_n - L \right| + \left| M - b_n \right| < \varepsilon
            \end{equation}

            We can apply the triangle inequality here.
            \begin{equation}
                \label{eq:3.22.b.10}
                \left| a_n - b_n - L + M \right| \leq \left| a_n - L \right| + \left| M - b_n \right| < \varepsilon
            \end{equation}

            Next, apply the inverse triangle inequality.
            \begin{equation}
                \label{eq:3.22.b.11}
                \left| \left| a_n - b_n \right| - \left| L - M \right| \right| \leq \left| a_n - b_n - L + M \right| < \varepsilon
            \end{equation}

            Recall the statement that $a_n \leq b_n$ and manipulate that.
            \begin{gather}
                \label{eq:3.22.b.12}
                a_n \leq b_n\\
                \label{eq:3.22.b.13}
                0 \leq b_n - a_n\\
                \label{eq:3.22.b.14}
                \left| b_n - a_n \right| = \left| a_n - b_n \right| = b_n - a_n
            \end{gather}

            Now, we can substitute in values for (\ref{eq:3.22.b.11}) elated to absolute values, specifically recalling (\ref{eq:3.22.b.7}) and (\ref{eq:3.22.b.14}).
            \begin{gather}
                \label{eq:3.22.b.15}
                \left| (b_n - a_n) - (L - M) \right| < \varepsilon
            \end{gather}

            This in turn implies that $b_n - a_n \to L - M$, which is a contradiction of (\ref{eq:3.22.b.3}).
            Ergo, $L \leq M$. %TENA.
        \end{proof}
\pagebreak
\setcounter{section}{31}

\section{Exercise 3.32}
    Let $(a_n) \subset \mathbb{R}$.
    Prove that if every subsequence of $(a_n)$ converges, then $(a_n)$ does.

    \subsection{Solution}
        \begin{proof}
            Suppose that we create a subsequence of $(a_n)$ that includes every member of $(a_n)$ except the first.
            By the rule, it converges.
            \begin{equation}
                \forall \varepsilon > 0, \exists N \in \mathbb{N} \backslash \{1\}: \forall n > N, \left| a_n - a \right| < \varepsilon
            \end{equation}

            Suppose that we do add $a_1$ to the sequence as its first member.
            There would be no violation of the rules of convergence.
            This would also be the sequence $(a_n)$ exactly.
            Therefore, $(a_n)$ converges.
        \end{proof}

\section{Exercise 3.33}

\pagebreak
\section{Exercise 3.34}
    Give an example of two sequences $(a_n)$ and $(b_n)$ that satisfy the following four conditions.
    \begin{multicols}{2}
        \begin{enumerate}[label=(\alph*)]
            \item   $\lim_{n \to \infty} a_n$ does not exist.
            \item   $\lim_{n \to \infty} b_n$ does not exist.
            \item   $(a_n \cdot b_n) \to \infty$
            \item   $(a_n \cdot b_{n+1}) \to -\infty$
        \end{enumerate}
    \end{multicols}

    \subsection{Solution}
        Suppose we define $(a_n)$ and $(b_n)$.
        \begin{gather}
            a_n := (-1)^n n\\
            b_n := (-1)^{n+1} n
        \end{gather}
    \pagebreak

\section{Exercise 3.36}
    Define a sequence $(a_n)$ recursively as follows. Let $a_1$ and $a_2$ be a pair of real numbers, and define recursively for all $n \geq 3$.
    \begin{equation*}
        a_n := \frac{a_{n - 1} + a_{n - 2}}{2}
    \end{equation*}
    Does $(a_n)$ necessarily converge?

    \subsection{Solution}
        Begin with $(a_n)$.
        \begin{equation}
            a_n := \frac{a_{n - 1} + a_{n - 2}}{2}
        \end{equation}

        By necessity, this is the average of the two and less than or equal to one of the prior values but greater than or equal to the other.
        
        Compute $\left| \delta a_n \right| = \left| a_n - a_{n-1} \right|$.
        \begin{equation}
            \left| \delta a_n \right|  =   \left| a_n - a_{n-1} \right| 
                =   \left| \frac{a_{n - 1} + a_{n - 2}}{2} - a_{n - 1} \right|
                =   \left| \frac{a_{n - 2} - a_{n - 1}}{2} \right|
        \end{equation}

        Let's check if $\left| \delta a_n \right|$ is monotone decreasing.
        Start with finding the value of $\delta a_{n+1}$.
        \begin{align}
            \delta a_{n+1}  &=   \frac{a_{n - 1} - a_n}{2}
                =   \frac{a_{n - 1} - \frac{a_{n - 1} + a_{n - 2}}{2}}{2}\\
                &=  \frac{2 a_{n - 1} - a_{n - 1} - a_{n - 2}}{4}
                =   \frac{a_{n - 1} - a_{n - 2}}{4}
        \end{align}

        This proof will involve some math, just to prove that $\left( \left| \delta a_n \right| \right)$ converges to 0, which will be done right now.
        It largely starts with a simple inequality of halves and absolute values.
        \begin{gather}
            \left| \frac{a_{n - 1} - a_{n - 2}}{2} \right| \leq \left| a_{n - 1} - a_{n - 2} \right|\\
            \left| \frac{a_{n - 1} - a_{n - 2}}{4} \right| \leq \left| \frac{a_{n - 1} - a_{n - 2}}{2} \right|\\
            \left| \delta a_{n + 1} \right| \leq \left| \delta a_n \right|
        \end{gather}

        $\left( \left| \delta a_n \right| \right)$ is monotone decreasing.
        We also can, from what we did here, develop an equation for $\left| \delta a_n \right|$ for $n \geq 3$.
        \begin{equation}
            \left| \delta a_n \right| = \left( \frac{1}{2} \right)^{n - 2} \left| \delta a_2 \right|
        \end{equation}

        Since $\left| \delta a_2 \right|$ is a constant and $\left( \frac{1}{2} \right)^{n - 2} \to 0$, then $\left| \delta a_n \right| \to 0$.
        \begin{gather}
            \forall \varepsilon > 0, \exists N : \forall n > N, \left| \delta a_n \right| < \varepsilon
        \end{gather}

        Substituting in the definition of $\left| \delta a_n \right|$ gives us something further.
        \begin{equation}
            \forall \varepsilon > 0, \exists N : \forall n > N, \left| a_n - a_{n - 1} \right| < \varepsilon
        \end{equation}

        We can use the triangle inequality to have this work for any $m, n > N$.
        \begin{gather}
            \left| a_n - a_{n - 1} \right| < \varepsilon\\
            \left| a_{n + 1} - a_n \right| < \varepsilon\\
            \left| a_{n + 1} - a_{n - 1} \right| \leq \left| a_{n + 1} - a_n \right| - \left| a_n - a_{n - 1} \right| < \varepsilon
        \end{gather}

        This can be expanded to work for any $m, n > N$.
        \begin{equation}
            \forall \varepsilon > 0, \exists N : \forall m, n > N, \left| a_m - a_n \right| < \varepsilon
        \end{equation}
        
        This in turn means that $(a_n)$ is Cauchy.
        This means that $(a_n)$ necessarily converges.

        % Compute $a_n$ for the first few $n$.
        % \begin{align}
        %     a_3 &=  \frac{a_2 + a_1}{2}\\
        %     a_4 &=  \frac{a_3 + a_2}{2}
        %         =   \frac{a_2 + a_1 + 2a_2}{4}
        %         =   \frac{3a_2 + a_1}{4}\\
        %     a_5 &=  \frac{a_4 + a_3}{2}
        %         =   \frac{1}{2} \left( \frac{3a_2 + a_1}{4} + \frac{a_2 + a_1}{2} \right)
        %         =   \frac{5a_2 + 3a_1}{8}\\
        %     a_6 &=  \frac{a_5 + a_4}{2}
        %         =   \frac{1}{2} \left( \frac{5a_2 + 3a_1}{8} + \frac{3a_2 + a_1}{4} \right)\\
        %         &=  \frac{5a_2 + 3a_1 + 6a_2 + 2a_1}{16}
        %         =   \frac{11a_2 + 5a_1}{16}
        % \end{align}
        
\section{Problem 3.38}
    Give an example of a monotone sequence that is not Cauchy.

    \subsection{Solution}
        Define $(a_n)$.
        \begin{equation}
            a_n := n
        \end{equation}

        It is monotonically increasing, but it diverges.
\pagebreak

\section{Problem 3.40}
    For each of the following, give an example of a sequence with that property or prove no such sequences exist.
    \begin{enumerate}[label=(\alph*)]
        \item   A Cauchy sequence with an unbounded subsequence.
        \item   An unbounded sequence with a Cauchy subsequence.
    \end{enumerate}

    \subsection{Solution (a)}
        \begin{proof}
            Recall the definition of a Cauchy sequence.
            \begin{equation}
                \forall\varepsilon > 0, \exists N : \forall m, n > N, \left| a_m - a_n \right| < \varepsilon
            \end{equation}

            We can rewrite the last part of this.
            \begin{gather}
                -\varepsilon < a_m - a_n < \varepsilon\\
                a_n - \varepsilon < a_m < a_n + \varepsilon
            \end{gather}

            This means that $a_m$ is bounded by the other possible $a_n$.
            There are values not bounded by this restriction, the ones where $m \leq N$.
            Suppose we found the maximum value of $a_m$ where $m \leq N$.
            \begin{equation}
                S = \left\{ a_1, a_2, \dots, a_N \right\}
            \end{equation}

            Suppose we add $a_n + \varepsilon$ and $a_n - \varepsilon$ to this and set $\varepsilon = 1$.
            \begin{multline}
                S   =   \left\{ a_1, a_2, \dots, a_N, a_n - \varepsilon, a_n + \varepsilon \right\}\\
                    =   \left\{ a_1, a_2, \dots, a_N, a_n - 1, a_n + 1 \right\}
            \end{multline}

            We can define a maximum $U$ and a minimum $L$ value of this, which set bounds for $(a_n)$.
            \begin{align}
                U &= \max(S)\\
                L &= \min(S)
            \end{align}

            As such, for all $m$, $L \leq a_m \leq U$.
            Suppose we call our subsequence $(a_{n_k})$.
            Members of $(a_{n_k})$ will all be members of $(a_n)$.
            This means that for all $n_k$, $L \leq a_{n_k} \leq U$.
            This means that it is bounded.
        \end{proof}

    \subsection{Solution (b)}
        I will define a sequence whose Cauchy subsequence is only the odd terms.
        \begin{equation}
            a_n := n^{(-1)^n}
        \end{equation}
\pagebreak

\section{Problem 3.42}
    Suppose the sequence $(x_n)$ is the sum of two other sequences, $(y_n)$ and $(z_n)$.
    That is $x_n = y_n + z_n$ for all $n$.
    If $(x_n)$ is bounded, and $(y_n)$ and $(z_n)$ are both monotone, must $(x_n)$ be convergent?
    What if $(y_n)$ and $(z_n)$ are also bounded?

    \subsection{Solution (a): Unbounded}
        In the first case, $(x_n)$ must not necessarily be convergent.
        Suppose we define two new sequences.
        \begin{gather}
            \delta y_n := y_{n + 1} - y_{n}\\
            \delta z_n := z_{n + 1} - z_{n}
        \end{gather}

        From here, we can define $\delta x_n$.
        \begin{equation}
            \delta x_n = \delta y_n + \delta z_n
                =   y_{n + 1} - y_{n} + z_{n + 1} - z_{n}
                =   x_{n + 1} - x_n
        \end{equation}

        For $(x_n)$ to converge, $(x_n)$ must be Cauchy.
        \begin{gather}
            \forall \varepsilon > 0, \exists N : \forall m, n > N, \abs{x_m - x_n} < \varepsilon
        \end{gather}

        It's possible that $m = n + 1$.
        \begin{gather}
            \abs{x_{n + 1} - x_n} = \abs{\delta x} 
                =   \abs{\delta y_n + \delta z_n}
                <   \varepsilon
        \end{gather}

        Now, suppose that $\delta y_n = \delta z_n = \frac{\varepsilon}{2}$.
        \begin{equation}
            \abs{\delta x}  =   \abs{\delta y_n + \delta z_n}
                =   \abs{\frac{\varepsilon}2 + \frac{\varepsilon}2}
                =   \abs{\varepsilon}
                =   \varepsilon
        \end{equation}

        This would not be possible if $(x_n)$ has to be convergent.
        Therefore, $(x_n)$ does not have to be convergent.

    \subsection{Solution (b): Bounded}
        Recall the monotone convergence theorem.
        It says that if a sequence is monotone, it converges if and only if it is bounded.
        This means that $(y_n)$ and $(z_n)$ are also Cauchy.
        \begin{gather} \label{eq:3.52}
            \forall \varepsilon > 0, \exists N_y : \forall m, n > N_y, \abs{y_m - y_n} < \varepsilon\\
            \label{eq:3.53}
            \forall \varepsilon > 0, \exists N_z : \forall m, n > N_z, \abs{z_m - z_n} < \varepsilon
        \end{gather}

        We can add the last terms of these two together, then use the Triangle Inequality.
        In this case, set $m, n > N = \max\{N_y, N_z\}$.
        \begin{align}
            &\forall \varepsilon > 0, \exists N_y : \forall m, n > N\\
            \varepsilon &>  \abs{y_m - y_n} + \abs{z_m - z_n}
                \geq    \abs{y_m - y_n + z_m - z_n}\\
            \varepsilon &>  \abs{y_m + z_m - y_n - z_n}
                =   \abs{x_m - x_n}
        \end{align}

        Put this together with (\ref{eq:3.52}) or (\ref{eq:3.53}).
        \begin{equation}
            \forall \varepsilon > 0, \exists N_y : \forall m, n > N, \abs{x_m - x_n} < \varepsilon
        \end{equation}

        This means that \seq{x} is Cauchy, which means that it converges.

    \subsection*{Scratch Work (b)}
        MCT says $y_n$ and $z_n$ both converge.
        \begin{gather*}
            \forall \varepsilon > 0, \exists N : \forall n > N, \abs{y_n - y} < \varepsilon\\
            \forall \varepsilon > 0, \exists N : \forall n > N, \abs{z_n - z} < \varepsilon\\
            \abs{y_n - y} + \abs{z_n - z} < 2\varepsilon\\
            \text{Cauchy: } 2\varepsilon > \abs{y_m - y_n} + \abs{z_m - z_n} \geq \abs{y_m - y_n + z_m - z_n}\\
            \abs{y_m - y_n + z_m - z_n} = \abs{y_m + z_m - y_n - z_n}
                =   \abs{x_m - x_n}
        \end{gather*}

        $(x_n)$ is Cauchy, so it converges.
\pagebreak

\section{Problem 3.44}
    Prove that if \seq{a} is a bounded sequence which does not converge, then it must contain two subsequences, both of which converge, but which converge to different values. 

    \subsection{Solution}
        \begin{proof}
            Bolzano-Weierstrass says every bounded sequence has a convergent subsequence.
            Naturally, \seq{a} also has one, which we could call \seq[n_k]{a}.
            Suppose we call the set of all $n_k$, $\mathbb{N}_k$.
            We can create three sets from these two sequences.
            \begin{gather}
                A = \{a_n | n \in \mathbb{N}\}\\
                A_k = \{a_{n_k} | n_k \in \mathbb{N}_k\}\\
                A_k^c = \{a_n | n \in \mathbb{N} \backslash \mathbb{N}_k\}
            \end{gather}

            Let $(m_j)_{j=1}^\infty$ be the increasing enumeration of $\mathbb N \setminus \mathbb N_k$.
            Then $(a_{m_j})_{j=1}^\infty$ is the sequence of terms of $(a_n)$ not contained in $(a_{n_k})$.

            For \seq{a} not to converge, \seq[m_j]{a} must be an infinite sequence.
            \begin{proof}[Proof by Contradiction]
                Suppose for contradiction that \seq[m_j]{a} is finite.
                That means that \seq[j]{m} is finite.
                Since it is finite, it has a maximum member, which we call $M$.
                For all $m > M$, $a_m \notin \seq[m_j]{a}$.
                By definition, for all $m > M$, $a_m \in \seq[n_k]{a}$.
                \seq[n_k]{a} converges, which we can use the definition of.
                \begin{equation}
                    \forall \varepsilon > 0, \exists N : \forall n_k > N, \abs{a_{n_k} - a} < \varepsilon
                \end{equation}

                For some values of $N$, $N < M$, so if $n > M$, then $n > N$.
                For others, $N > M$, so if $n > N$, then $n > M$.
                In either case, since \seq[m_j]{a} converges, we can rewrite its convergence definition to include $M$.
                \begin{equation}
                    \forall \varepsilon > 0, \exists N : \forall n > \begin{cases} N & \mathif N > M\\ M & \mathif M \geq N \end{cases}, \abs{a_n - a} < \varepsilon
                \end{equation}

                If we add back the values of \seq[m_j]{a}, this does not change. 
                Quite the opposite, it would say that \seq{a} converges. 
                This contradicts that \seq{a} does not converge. 
                Ergo, our intial assumption is wrong, and \seq[m_j]{a} is not finite.
            \end{proof}

            Past that, since \seq{a} is bounded, \seq[m_j]{a} is also bounded.
            This means that \seq[m_j]{a} either is convergent or has a convergent subsequence.
            If \seq[m_j]{a} converges to the same value as \seq[n_k]{a}, all subsequences of \seq{a} converge to the same number, which would mean that \seq{a} converges to that number, which means that \seq{a} converges, which we know cannot be the case.
            That would mean that \seq[m_j]{a} converges to something else.

            If \seq[m_j]{a} does not converge, it's possible that its convergent subsequence converges to the same value as \seq[n_k]{a}, but then we'd end up with another bounded, infinite subsequence.
            That could go on infinitely.
            If all of these converge to the same number, that means that all subsequences of \seq{a} converge to the same number, so \seq{a} converges, which is not the case.
            Ergo, one (or more) of them has to converge to a different number, which gives us two subsequences that converge to different limits.
        \end{proof}
\pagebreak

\section{Problem 3.45}
    Let \seq{a} be the sequence where $a_1 = 1$ and, for each $n > 1$, $a_n = a_{n-1} + \frac{1}{n^2}$.
    In 1734, Leonhard Euler famously proved that \seq{a} converges to $\frac{\pi^2}{6}$. Now let \seq{b} be the sequence where $b_1 = 1$ and, for each $n > 1$, $b_n = b_{n-1} + \frac{1}{n^3}$.
    Use Euler's result to prove that \seq{b} converges.

    \subsection{Solution}
        \begin{proof}
            Let's redefine $a_n$ and $b_n$.
            \begin{gather}
                a_n =   a_{n - 1} + \frac{1}{n^2}
                    =   1 + \sum_{i = 1}^{n} \frac{1}{i^2}
                    =   \sum_{i = 1}^{n} \frac{1}{i^2}\\
                b_n =   b_{n - 1} + \frac{1}{n^3}
                    =   1 + \sum_{i = 1}^{n} \frac{1}{i^3}
                    =   \sum_{i = 1}^{n} \frac{1}{i^3}
            \end{gather}

            Let's start with just the origin point of that.
            \begin{equation}
                i \geq 1
            \end{equation}

            Let's get closer to our definitions to $a_n$ and $b_n$.
            \begin{gather}
                i^3 \geq i^2\\
                \frac{1}{i^2} \geq \frac{1}{i^3}
            \end{gather}

            With every increment in $n$, $a_n$ increases by more than $b_n$.
            This means that for all $n > 1$, $a_n > b_n$.
            This gives \seq{b} an upper bound of \seq{a}, so \seq{a} is bounded.
            
            Furthermore, for all $i > 0$ (and by extension $n > 0$), $\frac{1}{i^3} > 0$, so \seq{b} is monotonically increasing.
            A monitonically increasing sequence bounded above converges.
            Ergo, \seq{b} converges.
        \end{proof}

    \subsection*{Scratch Work}
        \begin{gather*}
            a_n = 1 + \sum_{2}^{n} \frac{1}{i^2}\\
            b_n = 1 + \sum_{2}^{n} \frac{1}{i^3}\\
            \sum \frac{1}{i^2} \geq \sum\frac{1}{i^3}\\
            \frac{1}{i^2} \geq \frac{1}{i^3}
            i^3 \geq i^2\\
            i \geq 1
        \end{gather*}
\end{document}